\section{Conjunto de entrada branca}

Esta seção trata do conjunto \texttt{dados\_treino\_branco.csv}, em que a entrada é
aproximadamente branca. A Figura~\ref{fig:A_sinais} mostra os sinais de entrada e de
saída ao longo do tempo.

\begin{figure}[H]
\centering
\includegraphics[width=0.95\textwidth]{A_sinais.png}
\caption{Entrada $u[n]$ e saída $y[n]$ do conjunto de treino com entrada branca.}
\label{fig:A_sinais}
\end{figure}

\subsection{Média, variância e autocorrelação}

A Tabela~\ref{tab:A_stats} reúne as estatísticas de primeira e segunda ordem. A
entrada tem média próxima de zero e variância próxima de um, como se espera de um
ruído branco padrão. A saída tem média próxima de zero e variância menor, o que já
sugere um sistema que atenua parte da energia da entrada.

\begin{table}[H]
\centering
\caption{Estatísticas dos sinais do conjunto de entrada branca.}
\label{tab:A_stats}
\begin{tabular}{lcc}
\toprule
Sinal & Média & Variância \\
\midrule
$u[n]$ & \valAmediaU & \valAvarU \\
$y[n]$ & \valAmediaY & \valAvarY \\
\bottomrule
\end{tabular}
\end{table}

A Figura~\ref{fig:A_autocorr} apresenta as autocorrelações normalizadas. Para a
entrada, apenas uma fração de \valAfracAutoU{} dos atrasos diferentes de zero
ultrapassa a banda de confiança, valor próximo do esperado por acaso, o que é
compatível com um sinal branco. Para a saída, essa fração sobe para
\valAfracAutoY, o que indica memória: a saída é a versão filtrada da entrada e por
isso guarda correlação entre amostras vizinhas.

\begin{figure}[H]
\centering
\includegraphics[width=0.95\textwidth]{A_autocorr.png}
\caption{Autocorrelação normalizada da entrada e da saída (entrada branca). As linhas
tracejadas indicam a banda de confiança aproximada de 95\%.}
\label{fig:A_autocorr}
\end{figure}

\subsection{Correlação cruzada}

Para entrada branca vale a relação $R_{yu}[k] = \sigma_u^2\,h[k]$, ou seja, a
correlação cruzada reproduz a resposta ao impulso a menos de um fator de escala. A
Figura~\ref{fig:A_xcorr} mostra esse comportamento, com pico no atraso
$k = \valApicoXcorr$ e valores próximos de zero nos atrasos negativos, o que é
coerente com um sistema causal.

\begin{figure}[H]
\centering
\includegraphics[width=0.9\textwidth]{A_xcorr.png}
\caption{Correlação cruzada $R_{yu}[k]$ para o conjunto de entrada branca.}
\label{fig:A_xcorr}
\end{figure}

\subsection{Densidade espectral de potência}

A Figura~\ref{fig:A_psd} traz as densidades espectrais de potência da entrada e da
saída. O espectro da entrada é aproximadamente plano, o que confirma o caráter
branco do sinal. O espectro da saída concentra energia nas frequências baixas e cai
nas altas, indicando um sistema do tipo passa-baixas.

\begin{figure}[H]
\centering
\includegraphics[width=0.9\textwidth]{A_psd.png}
\caption{Densidade espectral de potência da entrada e da saída (entrada branca).}
\label{fig:A_psd}
\end{figure}

\subsection{Espectro cruzado}

A Figura~\ref{fig:A_espectro_cruzado} mostra o módulo e a fase do espectro cruzado
entre entrada e saída. O módulo segue a forma do espectro da saída, e a fase varia de
maneira regular nas frequências em que há energia, o que prepara a estimativa da
resposta em frequência.

\begin{figure}[H]
\centering
\includegraphics[width=0.9\textwidth]{A_espectro_cruzado.png}
\caption{Módulo e fase do espectro cruzado entrada--saída (entrada branca).}
\label{fig:A_espectro_cruzado}
\end{figure}

\subsection{Coerência}

A coerência, na Figura~\ref{fig:A_coerencia}, fica próxima de um nas frequências
baixas e cai nas altas. O valor médio é \valAcohMedia{} e a fração de frequências com
coerência acima de $0{,}9$ é \valAcohFrac. Essa queda não indica um problema: ela
ocorre justamente onde a saída tem pouca energia, de modo que a relação entrada--saída
é confiável na faixa de passagem do sistema e perde confiabilidade na faixa de
rejeição.

\begin{figure}[H]
\centering
\includegraphics[width=0.9\textwidth]{A_coerencia.png}
\caption{Coerência entre entrada e saída (entrada branca).}
\label{fig:A_coerencia}
\end{figure}

\subsection{Resposta em frequência}

A Figura~\ref{fig:A_resposta_frequencia} compara as duas estimativas da resposta em
frequência. Na faixa de passagem, em que a coerência é alta, os dois métodos
coincidem. Na faixa de rejeição, o primeiro método tende a superestimar o módulo,
porque o ruído presente na saída infla o termo $\hat{S}_{yy}$. O segundo método, que
usa o espectro cruzado, é menos sensível a esse ruído e fornece também a fase.

\begin{figure}[H]
\centering
\includegraphics[width=0.9\textwidth]{A_resposta_frequencia.png}
\caption{Resposta em frequência estimada pelos dois métodos (entrada branca).}
\label{fig:A_resposta_frequencia}
\end{figure}

\subsection{Identificação do modelo}

O modelo FIR foi ajustado por mínimos quadrados. Ao variar a ordem, o coeficiente de
determinação passa de \valArquadMcinco{} com cinco coeficientes para \valArquadMdez{}
com dez, e depois se estabiliza, o que sugere que a dinâmica relevante cabe em poucos
coeficientes. Adotou-se $M = 30$, com coeficiente de determinação de \valArquadTreino{}
no próprio conjunto de treino. A Figura~\ref{fig:A_fir_impulso} mostra a resposta ao
impulso estimada e a sua concordância com a correlação cruzada normalizada. O ganho
em frequência zero, dado pela soma dos coeficientes, é \valAganhoDC.

\begin{figure}[H]
\centering
\includegraphics[width=0.95\textwidth]{A_fir_impulso.png}
\caption{Resposta ao impulso do modelo FIR e comparação com a correlação cruzada
normalizada.}
\label{fig:A_fir_impulso}
\end{figure}

A Figura~\ref{fig:A_fir_ajuste} mostra um trecho da saída do conjunto de treino
sobreposta à previsão do modelo, que acompanha o sinal de perto.

\begin{figure}[H]
\centering
\includegraphics[width=0.9\textwidth]{A_fir_ajuste.png}
\caption{Saída do treino e previsão do modelo FIR (trecho inicial).}
\label{fig:A_fir_ajuste}
\end{figure}

A Figura~\ref{fig:A_fir_frequencia} compara a resposta em frequência do modelo FIR
com as estimativas não paramétricas. As três curvas coincidem na faixa de passagem,
o que dá confiança ao modelo identificado.

\begin{figure}[H]
\centering
\includegraphics[width=0.9\textwidth]{A_fir_frequencia.png}
\caption{Resposta em frequência do FIR identificado e das estimativas não
paramétricas.}
\label{fig:A_fir_frequencia}
\end{figure}

\subsection{Validação}

O modelo foi aplicado ao conjunto \texttt{dados\_validacao\_branco.csv}. A
Figura~\ref{fig:A_validacao} mostra um trecho da saída verdadeira e da saída prevista,
que se sobrepõem. O coeficiente de determinação na validação é \valArquadVal, o índice de
ajuste é \valAfitval\%, e o erro quadrático médio é \valArmseVal. Como esses valores
são próximos dos obtidos no treino, não há sinal de sobreajuste.

\begin{figure}[H]
\centering
\includegraphics[width=0.9\textwidth]{A_validacao.png}
\caption{Saída verdadeira e prevista no conjunto de validação (trecho inicial).}
\label{fig:A_validacao}
\end{figure}

A Figura~\ref{fig:A_residuos_validacao} mostra os resíduos da validação e a sua
autocorrelação. Os resíduos têm média próxima de zero e autocorrelação dentro da
banda de confiança, ou seja, são próximos de um ruído branco. Isso indica que o
modelo captura a parte previsível da saída e deixa pouco padrão por explicar.

\begin{figure}[H]
\centering
\includegraphics[width=0.95\textwidth]{A_residuos_validacao.png}
\caption{Resíduos da validação e a sua autocorrelação.}
\label{fig:A_residuos_validacao}
\end{figure}
